%!TEX root = ../../thesis.tex
\chapter{おわりに}
\label{chap:conclusion}

本論文では，目標方向の決定をシナリオが担い，各区間の走行を基盤モデルが担うシステムにより，走行対象の環境での事前走行や再学習，自己位置推定を必要とせずに，シナリオに従って目的地まで移動できるかを検証した．
シミュレータでの検証により，高層棟のデータのみでファインチューニングしたシステムが，学習環境の7例に加え，データを収集していない低層棟の6例においても，シナリオに従って目的地まで到達することを確認した．

本研究では，通路分類モジュールに\DINOvTwo と\UniDepth を，経路追従モジュールに\NoMaD を用いた．
いずれも特徴抽出器を凍結し，通路分類では3方向の有無を推定する部分を，経路追従では目標方向の埋め込みと拡散モデルを学習した．
これにより，用いるモデルのパラメータ数を合計約7500万に抑え，単眼RGBカメラのみで動作するシステムを構成した．

今後の課題を以下に示す．
\begin{itemize}
  \item ロボットに搭載した計算機を用いて，実環境で評価する．
  \item 基盤モデルの凍結や占有格子の併用など，各構成要素の効果を比較実験により定量的に示す．
  \item シナリオの数と試行回数を増やし，到達率や失敗の要因を評価する．
\end{itemize}

% TODO: 今後の課題を書き足す．
